\subsection*{What the perturbation data can support}
\label{sec:evaluability}

The preceding sections establish that features are biologically annotated, modularly organized and
causally specific. The question that remains is whether they encode regulatory logic: when a
transcription factor is knocked down, do the features that respond correspond to that factor's
known targets?

A recovery rate answers that question only against a reference. Before measuring the atlas we
therefore asked how much curated regulatory structure the assay can express at all. The CRISPRi
resource measures 6,546 genes. Of the 73 perturbed transcription factors with at least 50 cells,
only nine have five or more TRRUST targets among the measured genes, three have ten or more, and
one has twenty or more; the median perturbed factor has exactly \textbf{one} measurable target
(Fig.~\ref{fig:ceiling}b, \suptab{evaluability}). DoRothEA, whose regulons are larger, does
better---13 factors with five or more measured targets at the A+B+C confidence tier, median 15---but
the constraint is severe under every reference. For most of the panel, target-specific recovery is
not merely difficult; it is undefined, for any method.

\subsection*{A recoverability benchmark}
\label{sec:ceiling}

We therefore evaluated the SAE alongside methods that span the range of what is possible on these
data, each emitting a ranked list of 100 putative targets per factor and each scored by the same
hypergeometric enrichment with Benjamini--Hochberg correction across the panel and the measured
genes as the universe (Fig.~\ref{fig:ceiling}a, \suptab{ceiling}).

Differential expression of the knockdown itself is the perturbation-aware reference: it uses the
same cells and the actual intervention, and therefore bounds what any observational method can
achieve here. It recovers curated targets for 1/9 factors against TRRUST (11.1\%), 3/13 against
DoRothEA A+B+C (23.1\%) and 8/18 against DoRothEA at all confidence levels (44.4\%). Established
network inference does not approach this bound: GENIE3 reaches 12.5\%, 10.0\% and 0.0\% under the
three references, GRNBoost2 12.5\%, 10.0\% and 7.1\%, and simple co-expression 25.0\%, 20.0\% and
14.3\%. Contextual gene embeddings taken directly from the model, without any sparse decomposition,
reach 12.5\%, 10.0\% and 0.0\%. Size-matched random gene sets recover nothing under any reference,
confirming that the test is calibrated.

Two things follow. The ceiling on these data is low---even the intervention itself identifies
curated targets for fewer than half of the factors---and no observational method, whether a
purpose-built network-inference algorithm or a foundation-model representation, comes close to it.
The choice of universe matters for this calibration: scoring against a nominal 20,000-gene universe
rather than the 6,546 measured genes makes random gene sets appear significant for 16.7\% of
factors, because predictions can only be drawn from what is measured.

\subsection*{Perturbation responses are detectable but not transcription-factor specific}
\label{sec:perturbation}

We mapped perturbation responses at layer~11 using up to 200 CRISPRi cells per target against 600
non-targeting cells of the same line, encoding every cell through the atlas dictionary (Methods).
The panel comprises 20 evaluable transcription factors and, as a negative control, ten
perturbations of genes that are not transcription factors under any reference.

\paragraph{The unit of replication.} A perturbation is applied per cell, but a per-position
statistic treats each gene position as an observation. The intraclass correlation of feature
activation across positions within a cell is 0.000298, which at the observed 2,020 positions per
cell gives a design effect of 1.61: the effective sample size is about 62\% of the nominal position
count. The correction is worth making and we make it throughout---significance is assessed over
cells---but its magnitude is bounded, because TopK sparsity leaves activations at different
positions of one cell close to independent.

\paragraph{Responses are weak and non-specific.} Thresholding on effect size alone---a shift
exceeding half a control standard deviation, with no correction for testing 4,608 features---the
panel yields a mean of 3.20 responding features per factor, with 15 of 20 factors showing at least
one. Correcting for multiple testing across features collapses this to two of twenty, so most of
what a per-feature threshold registers does not survive the number of features being tested.
Assessed over cells at the same effect threshold, the mean is 1.90 responding features per factor (\suptab{cell_level}). The ten non-transcription-factor controls yield
\emph{more} responding features than the transcription factors (2.30 versus 1.90), so what the
features register is not specific to transcription-factor perturbation.

\paragraph{Where the sensitivity ends.} The clearest measurement of the limit compares the two
readouts on identical cells. Differential expression detects the knocked-down gene itself in 15 of
16 targets, at FDR $< 0.05$ with a median Cohen's $d$ of $-1.19$ (range $-2.76$ to $-0.49$). The
SAE feature response detects it in 1 of 20. The perturbation is therefore unambiguously present in
the expression profile of the same cells and is not registered by the feature responses. The
features are not insensitive because the data are thin; they track global cell-state shifts rather
than the gene-level consequences of the intervention. This is the specific sense in which detection
is limited, and it is a property of what the dictionary encodes.

\paragraph{Sample size.} Restricting to the 14 factors with at least 100 cells, so that sample size
is not confounded with which factors are available, the responding-feature count is flat across a
tenfold increase in cells (0.56 at $n{=}10$, 0.65 at $n{=}20$, 0.65 at $n{=}50$, 0.61 at
$n{=}100$). Detection of the knockdown rises from 5.0\% to 7.1\% and saturates. Nominal target
enrichment improves from $n{=}10$ to $n{=}50$ (7.7\% to 27.5\%) and then plateaus, so the earlier
choice of 20 cells was somewhat under-powered, but the gain saturates well before the effect becomes
substantial, and under correction across the panel no sample size yields a significant factor
(\suptab{sweep}).

\paragraph{Widening the training distribution.} A dictionary trained on the same control cells
pooled with 3,000 Tabula Sapiens cells from three primary tissues gives 2.55 responding features per
factor against 1.90 for the atlas dictionary, with the same non-transcription-factor controls at
2.30 and detection of the knockdown at 0 of 20. Diversifying the training distribution beyond
immortalized lines does not make regulatory structure recoverable.

\subsection*{Replication across cell lines and independent datasets}
\label{sec:replication}

The measurement was repeated in three further cell lines from the same resource, with perturbed and
non-targeting cells drawn from the same line in every case, and on two perturbation datasets
generated by other groups with other protocols and substantially larger gene panels
(Fig.~\ref{fig:power}c, \suptab{crossline}).

Across K562, RPE1, Jurkat and HepG2 the pattern is the same. Pooling the four lines,
\textbf{1 of 66} evaluable transcription factors shows target enrichment at FDR $< 0.05$
(1.5\%): none of 18 in K562, none of 16 in RPE1, one of 16 in Jurkat and none of 16 in HepG2. The
knocked-down gene itself is detected in 2 of 68 factors across the four lines. Responding-feature
counts are low throughout (means of 0.06--1.90 per factor) and do not separate transcription
factors from the non-transcription-factor controls, which average 0.17--2.30 in the same lines. The
design effect is 1.61--1.73 everywhere, so the correction for the unit of replication behaves
consistently across backgrounds. RPE1 is a non-transformed retinal pigment epithelial line and
HepG2 a hepatocyte line, so the result is not confined to the leukaemic backgrounds in which these
models are most often probed.

Both external datasets measure far more of the transcriptome than the primary resource, and the
difference changes the result. In THP-1 monocytes 18,649 genes are measured and the median
perturbed factor has 19.5 curated TRRUST targets among them; in the K562 Perturb-seq data 33,694
genes are measured. The Perturb-seq screen is a CRISPR \emph{activation} experiment, so it also
tests the question under gene induction rather than repression: the median perturbed factor is
induced with a Cohen's $d$ of $+0.98$ and yields 405 differentially expressed genes.

On these data the features do recover targets. Against DoRothEA the SAE reaches 20--22\% of factors
in the THP-1 data and 23--25\% in the Perturb-seq data, where on the primary panel it reached none.
Against TRRUST it reaches none in either, while differential expression on the same cells reaches
18\% and 30\%.

Whether that recovery is transcription-factor-specific depends on the reference, and a control
settles which part of it to believe. We ran ten perturbations of genes absent from every reference
network through the identical pipeline and tested their predicted sets against the \emph{same}
regulons as the transcription factors. Against the high-confidence DoRothEA subset, a factor's own
predicted set is enriched for its own regulon in 2 of 13 cases (15\%) while the
non-transcription-factor predictions are enriched for those same regulons in 1 of 117 (1\%) --- a
fifteenfold separation, so this signal is specific. Against the union of all confidence levels the
separation largely disappears: 7 of 16 (44\%) for the matching factors against 39 of 144 (27\%) for
the controls. Those regulons have a median of 402 measured genes and reach 9,321, and a broad
expression signature enriches for them regardless of which gene was perturbed. The
all-confidence numbers therefore do not measure regulatory recovery, and we do not read them as
such.

Two conclusions follow, and they pull in different directions. The absence of detectable
transcription-factor specificity in the primary CRISPRi panel is substantially a property of that
assay: with a median of one measurable target per factor, neither these features nor differential
expression on the same cells nor established network inference can demonstrate specificity there.
Where the transcriptome is measured and the perturbation is strong, the features carry a modest but
genuinely factor-specific target signal. Fifteen percent of factors, against a one percent control
rate, is a small fraction of a small panel; it is also clearly not zero.
